
Every additional nat of KL divergence applied during RLHF fine-tuning increases the gap between a reward model's evaluation score and actual human preference by approximately 0.143 units on a normalized scale --- a relationship consistent enough to explain 96\% of divergence variance across 10 optimization checkpoints and to replicate with slopes within 12\% across two independent model families. The more aggressively a model is optimized for standard alignment metrics, the less calibrated that model becomes to held-out human judgment.

This finding runs counter to the operating assumption of RLHF. Reinforcement Learning from Human Feedback --- the dominant paradigm for making large language models helpful, harmless, and honest --- promises that optimizing against a reward model trained on human preferences should make a model \textit{more} aligned with human judgment. Yet the data reveal a systematic divergence: reward model scores climb monotonically with KL budget, while gold human preference peaks early (at approximately 2 nats of KL divergence in the Coste et al. dataset) and subsequently declines by 40\%. The more optimized the model, the wider the gap between what the evaluation instrument reports and what held-out human judges prefer.

\textbf{The surface problem} is well-established: RLHF trains against a proxy metric --- the reward model score --- that approximates but is not equivalent to gold human preference. As optimization pressure increases, the policy finds outputs that satisfy the proxy while diverging from the underlying target. This ''reward hacking'' or ''overoptimization'' phenomenon has been described qualitatively in the literature [Coste et al., 2023; Gao et al., 2023]: RM scores rise while human preference metrics plateau or reverse.

\textbf{The deeper problem} is that this divergence has not been quantitatively characterized as a function of optimization pressure. Prior work does not report \textit{how fast} the gap grows, whether the growth rate is model-family-specific, or whether a standardized measurement instrument can enable cross-study comparison. Without these quantitative answers, practitioners cannot estimate when RLHF optimization degrades calibration past an acceptable threshold, nor can researchers compare overoptimization severity across experimental settings.

\textbf{The gap this work addresses} lies at the intersection of three literatures that have not been connected: RLHF overoptimization theory, bidirectional human-AI alignment measurement, and Goodhart's Law in machine learning. Each literature describes part of the phenomenon, but none provides a named construct, a quantified slope, or a normalized metric enabling cross-dataset comparison. The ICLR 2025 Workshop on Bidirectional Human-AI Alignment synthesized 400 papers and found that existing alignment benchmarks systematically measure only the $AI\rightarrow Human$ direction --- reward model scores, safety classifiers, benchmark accuracy --- while $Human\rightarrow AI$ calibration (whether held-out human preference tracks model optimization trajectories) is studied in a parallel HCI track not integrated with ML evaluation. RLHF overoptimization is the empirical instantiation of this bidirectional alignment tension, where optimizing $AI\rightarrow Human$ proxy metrics systematically degrades evaluation calibration to held-out human judgment.

\textbf{The key insight} is that the normalized divergence gap --- RM\_norm $-$ gold\_preference, where both signals are min-max scaled to [0,1] --- grows linearly with KL budget at a consistent, measurable rate. This normalized metric removes scale artifacts while preserving direction and monotonicity, creating a measurement instrument that functions across different experimental settings. Applying ordinary least squares (OLS) regression to this metric across published RLHF experimental data reveals a named quantity --- the \textit{calibration-alignment divergence curve} --- characterized by a slope $\beta$ that is significantly positive, statistically robust, and reproducible across independent model families.

This work makes four contributions:

\begin{enumerate}
\item \textbf{The calibration-alignment divergence construct.} We define and operationalize the calibration-alignment divergence gap as a named bidirectional alignment construct --- the first regression characterization of the normalized divergence gap (RM\_norm $-$ gold\_preference) as a function of KL optimization budget with cross-dataset slope comparison. Prior work described the phenomenon qualitatively; we compute $\beta, R^2$, parametric confidence intervals, and bootstrap confidence intervals, enabling systematic cross-study comparison.
\end{enumerate}

\begin{enumerate}
\item \textbf{Quantified divergence curve slopes.} In Coste et al. [2023] data, $\beta$ = 0.1433 $nat^{-1}$ ($R^2$ = 0.9577, p = 8.89 $\times 10^{-7}$; t = 13.461, F = 181.2, N = 10) --- a near-perfect linear fit. The calibration-alignment divergence gap grows by approximately 0.143 normalized units per nat of KL budget, a relationship that explains 96\% of gap variance across 10 optimization checkpoints.
\end{enumerate}

\begin{enumerate}
\item \textbf{Cross-dataset replication with effect size comparison.} In independent Gao et al. [2023] data (different model family, different scale, 6B reward model), $\beta$ = 0.1599 $nat^{-1}$ (p = 0.0025, $R^2$ = 0.7008, N = 10), yielding a cross-dataset slope ratio of 1.116. Two entirely independent experimental settings produce slopes within 12\% of each other --- consistent evidence (n = 2 datasets; further replication required to establish universality) that this pattern is not a model-family artifact.
\end{enumerate}

\begin{enumerate}
\item \textbf{The normalized divergence gap metric.} RM\_norm $-$ gold\_preference $\in$ [$-1$, +1] as a standardized evaluation calibration divergence instrument enabling cross-dataset comparison by removing raw scale differences between experimental settings.
\end{enumerate}

Section 2 situates this work in three strands of prior research. Section 3 describes the methodology and design rationale for the normalized gap metric. Section 4 presents the experimental setup and the four-step mechanistic verification. Section 5 reports results. Section 6 discusses implications, limitations, and connections to the broader bidirectional alignment agenda. Section 7 concludes.

